\documentclass[12pt]{article}
\usepackage[utf8]{inputenc}
\usepackage[T1]{fontenc}
\usepackage{hyperref}
\title{Torwali Language Resources for K-12 Curriculum Development}
\author{Compiled by Claude Sonnet 5 (claude-sonnet-5) for Andr\'as Kornai}
\date{}
\begin{document}
\maketitle
\section{Introduction}
Torwali (ISO 639-3 trw) is an Indo-Aryan language of Swat, Pakistan. It is endangered, and estimates of the number of speakers range from about 80,000 to 120,000.
\section{Primary Educational Resources}
IBT (Idara Baraye Taleem-o-Taraqi), founded by Zubair Torwali, promotes mother-tongue-based multilingual education in Swat. Its schools begin instruction in Torwali before moving students to Urdu and English. According to Wikipedia, IBT developed a Torwali writing system between 2005 and 2008 and produced children's textbooks, an alphabet book and a primer. These items were not located individually.
\begin{itemize}
\item \url{https://torwali.omeka.net/exhibits/show/torwaliexhibit/community-response}
\item \url{https://endangeredlanguages.com/node/112285}
\item \url{https://en.wikipedia.org/wiki/Torwali_language}
\end{itemize}
\section{Linguistic and Reference Materials}
\begin{itemize}
\item Zubair Torwali, ZAS Papers in Linguistics 64 (2020). It covers the orthography, the IBT schools, and adapting the MAIN narrative assessment tool to Torwali: \url{https://publikationen.ub.uni-frankfurt.de/opus4/files/55890/ZASPiL_64_30_Torwali.pdf}
\item Writings by Zubair Torwali: \url{https://criterion-quarterly.com/?p=6811} and \url{https://criterion-quarterly.com/?p=6529}
\end{itemize}
\section{Religious and Community Resources}
IBT is the community organisation (see the exhibit page above). No scripture was found in the search results.
\section{Digital and Audio Resources}
The Torwali character set has Unicode designations, according to Wikipedia. A sample entry is at \url{https://endangeredlanguages.com/lang/3501/samples/11346?hl=fr}. No audio or corpora were found.
\section{Curriculum Development Context}
Torwali is not part of Pakistan's public education system, according to the Endangered Languages Project. The IBT mother-tongue programme is the main existing model. Current enrolment and funding are unknown.
\end{document}